\input{../tex-inputs/latex-leseansicht-vorspann}

\section[Gustav Schwarzkopf an Arthur Schnitzler, 4. 8. 1900]{L04302 Gustav Schwarzkopf an Arthur Schnitzler, 4. 8. 1900}
\nopagebreak\mylabel{L04302v}
\rehead{ }\normalsize\beginnumbering\briefempfaengerindex{Schnitzler, Arthur@\textsc{Schnitzler, Arthur}!zzzSchwarzkopf, Gustav@\emph{von Gustav Schwarzkopf}!1900-08-042@{4. 8. 1900}|(be}
\toendnotes[C]{\smallbreak\pagebreak[2]}
\correspDesc{Versand  durch Gustav Schwarzkopf am 4. 8. 1900 in Wien
\newline{}Erhalt  durch Arthur Schnitzler im Zeitraum [5. 8. 1900 – 9. 8. 1900?] in Bad Ischl}\toendnotes[C]{\smallbreak}
\Standort{CUL, Schnitzler, B 96a.}
\physDesc{Brief, 1 Blatt, 4 Seiten, 1753 Zeichen
\newline{}Handschrift: schwarze Tinte, deutsche Kurrent}\toendnotes[C]{\smallbreak}
\pstart
           \raggedleft{}{\pb}4. VIII. 0\textcolor{gray}{0}\pend
           \vspace{0.5em}
\pstart
           Lieber Arthur, es wird{ }ſich leider nicht machen laſſen. Von der
               großen Fußtour mit Bahn und Fuß hätte ja für mich nie die Rede{ }ſein kö{\geminationn}en, aber einige Tage Salzburg\oindex{Salzburg@\textbf{Salzburg}, \emph{Verwaltungsgebiet}|pw} wäre ja zu leiſten geweſen. Nun tritt aber meine Köchin\pwindex{?? [Köchin von Gustav und Max Schwarzkopf, 1900] @\textsc{?? [Köchin von Gustav und Max Schwarzkopf, 1900]}|pwv}, wie ich Ihnen{ }ſchon mitgeteilt, eine
               Pilgerfahrt an, und zwar von 12.{[}{]}{[},{]} wie es jetzt feſtgeſetzt iſt.\footnote{\noindent{}und ko{\geminationm}t am 17. zurück.} Urlaub hab ich ihr{ }ſchon früher bewilligt, ich ka{\geminationn} ihn jetzt nicht mehr zurücknehmen ohne mich mit ihr zu verfeinden;{ }ſie wäre fähig
               die äußerſten Konſequenzen zu ziehen, we{\geminationn} ich{ }ſie um die
               Möglichkeit {\pb}bringe, ihre Sünden
               abzubüßen, und einer{ }ſolchen Gefahr will ich mich doch nicht ausſetzen. Ich ka{\geminationn} aber Max\pwindex{Schwarzkopf, Max 12.\,6.\,1857 Wien – 14.\,4.\,1928 ebd.@\textsc{Schwarzkopf, Max} (12.\,6.\,1857 Wien – 14.\,4.\,1928 ebd.), \emph{Rechtsanwalt}|pw} in der
               Wohnung\oindex{Wien@\textbf{Wien}!I., Innere Stadt@\textbf{I., Innere Stadt}!Tiefer Graben 23@\textbf{Tiefer Graben 23}, \emph{Wohngebäude}|pwv} nicht ganz allein laſſen,
               ich hätte da{\geminationn} wirklich keine ruhige Stunde in Salzburg\oindex{Salzburg@\textbf{Salzburg}, \emph{Verwaltungsgebiet}|pw}. Es iſt alſo nicht zu machen. Ich werde
               alſo wahrſcheinlich den ganzen Auguſt in Wien\oindex{Wien@\textbf{Wien}, \emph{Verwaltungsgebiet}|pw}
               bleiben, wie im vergangnen Jahr. Wie{ }ſich das Wetter jetzt anläßt, dürfte es ganz gut
               zu ertragen{ }ſein. Sie{ }ſchicken mir von den verſchiedenen Stationen Anſichtskarten und
                  we{\geminationn} Sie im September nach Wien\oindex{Wien@\textbf{Wien}, \emph{Verwaltungsgebiet}|pw} ko{\geminationm}en, werden Sie mir ja
               erzählen, was ich verloren habe, wie lange {\pb}Sie bei den verschiedenen Trödlern
               waren und wie oft die Geſellſchaft\pwindex{Beer-Hofmann, Richard 11.\,7.\,1866 Wien – 26.\,9.\,1945 New York City@\textsc{Beer-Hofmann, Richard} (11.\,7.\,1866 Wien – 26.\,9.\,1945 New York City), \emph{Schriftsteller}|pwv}\pwindex{Van-Jung, Leo 15.\,10.\,1866 Odessa – 2.\,7.\,1939 Riga@\textsc{Van-Jung, Leo} (15.\,10.\,1866 Odessa – 2.\,7.\,1939 Riga), \emph{Gesangspädagoge, Mathematiker}|pwv}\pwindex{Kerr, Alfred 25.\,12.\,1867 Breslau – 12.\,10.\,1948 Hamburg@\textsc{Kerr, Alfred} (25.\,12.\,1867 Breslau – 12.\,10.\,1948 Hamburg), \emph{Schriftsteller, Kritiker}|pwv}\pwindex{Goldmann, Paul 31.\,1.\,1865 Breslau – 25.\,9.\,1935 Wien@\textsc{Goldmann, Paul} (31.\,1.\,1865 Breslau – 25.\,9.\,1935 Wien), \emph{Schriftsteller, Journalist}|pwv} nicht – ganz einig
               war. Grüßen Sie übrigens alle Teilnehmer\pwindex{Beer-Hofmann, Richard 11.\,7.\,1866 Wien – 26.\,9.\,1945 New York City@\textsc{Beer-Hofmann, Richard} (11.\,7.\,1866 Wien – 26.\,9.\,1945 New York City), \emph{Schriftsteller}|pwv}\pwindex{Van-Jung, Leo 15.\,10.\,1866 Odessa – 2.\,7.\,1939 Riga@\textsc{Van-Jung, Leo} (15.\,10.\,1866 Odessa – 2.\,7.\,1939 Riga), \emph{Gesangspädagoge, Mathematiker}|pwv}\pwindex{Kerr, Alfred 25.\,12.\,1867 Breslau – 12.\,10.\,1948 Hamburg@\textsc{Kerr, Alfred} (25.\,12.\,1867 Breslau – 12.\,10.\,1948 Hamburg), \emph{Schriftsteller, Kritiker}|pwv}\pwindex{Goldmann, Paul 31.\,1.\,1865 Breslau – 25.\,9.\,1935 Wien@\textsc{Goldmann, Paul} (31.\,1.\,1865 Breslau – 25.\,9.\,1935 Wien), \emph{Schriftsteller, Journalist}|pwv}, besondere \textsc{Beer Hofma{\geminationn}\pwindex{Beer-Hofmann, Richard 11.\,7.\,1866 Wien – 26.\,9.\,1945 New York City@\textsc{Beer-Hofmann, Richard} (11.\,7.\,1866 Wien – 26.\,9.\,1945 New York City), \emph{Schriftsteller}|pw}} und Goldma{\geminationn}\pwindex{Goldmann, Paul 31.\,1.\,1865 Breslau – 25.\,9.\,1935 Wien@\textsc{Goldmann, Paul} (31.\,1.\,1865 Breslau – 25.\,9.\,1935 Wien), \emph{Schriftsteller, Journalist}|pw}. Es{ }ſind
               übrigens noch i{\geminationm}er Leute in Wien\oindex{Wien@\textbf{Wien}, \emph{Verwaltungsgebiet}|pw}, bei der letzten Vorſtellung\eventindex{Theater in der Josefstadt@\textbf{Theater in der Josefstadt}!Aufführung von Brand, 31.7.1900@Aufführung von Brand, 31.7.1900|pwv}
               der \textsc{Seccessions}-Bühne\orgindex{Secessionsbühne@Secessionsbühne|pw} »Brand\pwindex{\textcolor{red}{\textsuperscript{XXXX indx1}}!Brand. Ein dramatisches Gedicht in 5 Aufzügen (9 Bildern)@\strich\emph{Brand. Ein dramatisches Gedicht in 5 Aufzügen (9 Bildern)}|pw}« war es wieder{ }ſehr gut beſucht und man{ }ſah{ }ſogar Leute, denen es zur
               freundlichen Gewohnheit geworden iſt, gena{\geminationn}t zu werden.–
               Eben ko{\geminationm}e ich vom Weſtbahnhof\oindex{Wien@\textbf{Wien}!XV., Rudolfsheim-Fünfhaus@\textbf{XV., Rudolfsheim-Fünfhaus}!Westbahnhof@\textbf{Westbahnhof}, \emph{Bahnhof}|pw}. Ich habe Emil\pwindex{Schwarzkopf, Emil 17.\,9.\,1851 Wien – 28.\,1.\,1928 ebd.@\textsc{Schwarzkopf, Emil} (17.\,9.\,1851 Wien – 28.\,1.\,1928 ebd.), \emph{Übersetzer, Komponist}|pw}, dem kühnen Pariſer\oindex{Paris@\textbf{Paris}, \emph{Hauptstadt}|pw}-Reiſenden das Geleit gegeben; (er wäre
               unglücklich geweſen, we{\geminationn} ich das unterlaſſen hätte.)–
                  {\pb}Empfehlen Sie mich, bitte, Ihrer
                  Mama\pwindex{Schnitzler, Louise 8.\,7.\,1840 Kőszeg – 9.\,9.\,1911 Wien@\textsc{Schnitzler, Louise} (8.\,7.\,1840 Kőszeg – 9.\,9.\,1911 Wien)|pwv} und grüßen Sie alle
               Perſonen in \textsc{Ischl\oindex{Bad Ischl@\textbf{Bad Ischl}|pw}}, die ein Anrecht haben, von mir Grüße zu verlangen.\pend
           
\pstart
           Herzlichſt{\\[\baselineskip]}Ihr{\\[\baselineskip]}\spacefill\mbox{Gustav.}\pend
           \leftskip=0em{}\selectlanguage{ngerman}\endnumbering\briefempfaengerindex{Schnitzler, Arthur@\textsc{Schnitzler, Arthur}!zzzSchwarzkopf, Gustav@\emph{von Gustav Schwarzkopf}!1900-08-042@{4. 8. 1900}|)be}\mylabel{L04302h}
\begin{anhang}
\end{anhang}\newcommand{\dateiname}{L04302}\newcommand{\titel}{Gustav Schwarzkopf an Arthur Schnitzler, 4. 8. 1900}\newcommand{\editorInnen}{Herausgegeben von Jahnke, SelmaMüller, Martin Anton}\input{../tex-inputs/latex-leseansicht-abspann}
